\chapter{Introduction}

\section{Background}

Shortly taking on theorethical background, diffusion probabilistic models specify (i) forward Markov chain that, by repeatedly addng small amounts of Gaussian noise, converts an observed sample
$x_0$ into a near-isotropic Gaussian $x_T$, and
(ii) learned reverse-time chain $p_\theta(x_{t-1}\!\mid x_t)$ that removes noise step by step to
recover samples from data distribution \cite{sohl2015noneq,ho2020ddpm}.

In image generation \cite{ho2020ddpm},
this training recipe couples a noise schedule with a denoising network. Sampling becomes an
iterative sequence of steps of condtional reconstructions from $x_T\!\sim\!\mathcal{N}(0,I)$ to $x_0$. 

These components transfer well to time series forecasting by treating unknown future segment as a
conditional generation problem: given observed history -- conditional context window, a diffusion model can sample entire future
trajectories by denoising noise variables.

Feature-centric diffusion TSF keeps standard Gaussian coruption and places capacity in the
conditioning path. With future block $\mathbf{x}_0\in\mathbb{R}^{H\times D}$ and context
$\mathbf{x}_{-L+1:0}\in\mathbb{R}^{L\times D}$,
\[
\mathbf{x}_t=\sqrt{\bar\alpha_t}\,\mathbf{x}_0+\sqrt{1-\bar\alpha_t}\,\boldsymbol{\epsilon},\quad
\boldsymbol{\epsilon}\sim\mathcal{N}(\mathbf{0},\mathbf{I}),
\]
and denoiser predicts noise under context features
\[
\hat{\boldsymbol{\epsilon}}=\epsilon_\theta(\mathbf{x}_t,\mathbf{c},t),\qquad
\mathbf{c}=f_\phi(\mathbf{x}_{-L+1:0}),
\]
so reverse sampler uses $\hat{\boldsymbol{\epsilon}}$ to update $\mathbf{x}_{t-1}$ at each
noise level. Architectural decisions therefore target $f_\phi$, fusion operator between
$\mathbf{c}$ and the denoiser (e.g., channel-wise concat, cross-attention, FiLM-style modulation).

Diffusion TSF inherits discretized reverse-time sampler. Practical inference commonly requires
$\mathcal{O}(10^2\text{--}10^3)$ \cite{hu2024flowts} drift/denoiser evaluations per sample, making horizon-level
genration expensive even when the conditioner is strong. Flow Matching \cite{lipman2022flowmatching} replaces the discrete reverse chain with a continuous-time transport ODE,
\[
\frac{d\mathbf{z}(t)}{dt}=\mathbf{v}_\theta(\mathbf{z}(t),t,\mathbf{c}),\qquad
\mathbf{c}=f_\phi(\mathbf{x}_{-L+1:0}),
\]
and trains $\mathbf{v}_\theta$ on direct vector-field regression on a chosen probability path,
without simulating an SDE/Markov chain.
This targets the main computational bottleneck: sampling becomes numerical ODE integration with
coarse step sizes when learned trajectories are near-straight rather than hundreds of
noise-removal iterations. 

Rectified flow \cite{liu2022rectified} explicitly optimizes for straight transport trajectories, it
yields accurate generation under coarse (even single-step) discretization in inference phase.

\section{Literature Review}

Early work, such as TimeGrad \cite{rasul2021timegrad}, couples diffusion with an RNN encoder and autoregressive decoding scheme. RNN compresses the prefix into hidden states and the model samples future steps sequentially.
CSDI \cite{tashiro2021csdi} replaces recurrence with transformer self-attention. This enables direct time-index interactions in
the conditioning representation. This way of modelling directly supports blockwise recovery of $\hat{\mathbf{x}}_{1:H}$ in a
single reverse trajectory.

TimeDiff \cite{shen2023timediff} alters how target-side structure enters the conditional score. Future mixup injects partiallly revealed target components during training. Autoregressive initilization seeds the reverse chain with an explicit AR forecast to shift early
iterations toward a residual refinement.

Beyond feature encoders, S2DBM \cite{yang2024s2dbm} introduces a Brownian-Bridge series-to-series diffusion process with
pinned endpoints. While ARMD \cite{gao2024armd} swaps endpoint semantics, treating the future block as the diffusion
start state and the historical block as the terminal state, training reverse dynamics along this
trajectory.

FMTS \cite{hu2025fmts} instantiates rectified flow / flow matching on time-series blocks by transporting an entire
future segment while conditioning on a context encoder, avoiding diffusion-style noise schedules and
iterative denoising loops.
FlowTS \cite{hu2024flowts} further develops this line with a rectified-flow ODE backbone for time-series generation and
motivates the approach as a compute-efficient alternative to diffusion sampling in high-dimensional
sequence spaces.

\subsection{Stationarity Concerns}

Major concern few paper actually discuss and try to mitigate is non-stationarity of the time-series data they work with. Generative time-series models existing prioritize architectural priors over manual differencing to handle non-stationarity by simply taking first-differences -- common practice in time-series econometrics.

Distribution shift is primarly mitigated using Reversible Instance Normalization (RevIN) \cite{kim2021revin, shen2023timediff} or series stationarization backbones \cite{liu2022nonstationary} that restore temporal statistics during inference. Some frameworks, such as \cite{wang2024falda, mrdiff, yuan2024diffusionts, cdpm}, leverage seasonal-trend decomposition to isolate determinstic trends from stochastic components. FALDA \cite{wang2024falda} utilizes an NS-Adapter to mitigate epistemic uncertainty. LDT \cite{feng2024ldt} implements statistics-aware autoencoders with adaptive variance updating layers to stabilize latent space training against high-variance deviations.

\section{Contribution}

Inspired by the advancements in generative modelling in form of state of the art method -- Flow Matching, as well as opportunities in applying financial econometrics principles and methods of modelling, I propose FITS -- framework for time-series conditional probabilistic imputation in form of forecasting with decomposition of interpretable time-series components.

Main contributions are as follows:

\begin{enumerate}
    \item Application of the simpliest, yet by far the most efficient method of working with non-stationary time-series -- taking of the first differences.
    \item The broadening of interpretation capabilities by explicitely utilizing jump component within DiffusionTS \cite{yuan2024diffusionts} architecture.
    \item The usage of SoTA Flow Matching within conditional task of time-series forecasting.
\end{enumerate}
